% Options for packages loaded elsewhere
% Options for packages loaded elsewhere
\PassOptionsToPackage{unicode}{hyperref}
\PassOptionsToPackage{hyphens}{url}
\PassOptionsToPackage{dvipsnames,svgnames,x11names}{xcolor}
%
\documentclass[
  11pt,
  a4paper,
]{ctexrep}
\usepackage{xcolor}
\usepackage[top=2.5cm,bottom=2.5cm,left=2.5cm,right=2.5cm]{geometry}
\usepackage{amsmath,amssymb}
\setcounter{secnumdepth}{5}
\usepackage{iftex}
\ifPDFTeX
  \usepackage[T1]{fontenc}
  \usepackage[utf8]{inputenc}
  \usepackage{textcomp} % provide euro and other symbols
\else % if luatex or xetex
  \usepackage{unicode-math} % this also loads fontspec
  \defaultfontfeatures{Scale=MatchLowercase}
  \defaultfontfeatures[\rmfamily]{Ligatures=TeX,Scale=1}
\fi
\usepackage{lmodern}
\ifPDFTeX\else
  % xetex/luatex font selection
  \setmainfont[]{Liberation Serif}
  \setsansfont[]{Liberation Sans}
  \setmonofont[]{Liberation Mono}
\fi
% Use upquote if available, for straight quotes in verbatim environments
\IfFileExists{upquote.sty}{\usepackage{upquote}}{}
\IfFileExists{microtype.sty}{% use microtype if available
  \usepackage[]{microtype}
  \UseMicrotypeSet[protrusion]{basicmath} % disable protrusion for tt fonts
}{}
\makeatletter
\@ifundefined{KOMAClassName}{% if non-KOMA class
  \IfFileExists{parskip.sty}{%
    \usepackage{parskip}
  }{% else
    \setlength{\parindent}{0pt}
    \setlength{\parskip}{6pt plus 2pt minus 1pt}}
}{% if KOMA class
  \KOMAoptions{parskip=half}}
\makeatother
% Make \paragraph and \subparagraph free-standing
\makeatletter
\ifx\paragraph\undefined\else
  \let\oldparagraph\paragraph
  \renewcommand{\paragraph}{
    \@ifstar
      \xxxParagraphStar
      \xxxParagraphNoStar
  }
  \newcommand{\xxxParagraphStar}[1]{\oldparagraph*{#1}\mbox{}}
  \newcommand{\xxxParagraphNoStar}[1]{\oldparagraph{#1}\mbox{}}
\fi
\ifx\subparagraph\undefined\else
  \let\oldsubparagraph\subparagraph
  \renewcommand{\subparagraph}{
    \@ifstar
      \xxxSubParagraphStar
      \xxxSubParagraphNoStar
  }
  \newcommand{\xxxSubParagraphStar}[1]{\oldsubparagraph*{#1}\mbox{}}
  \newcommand{\xxxSubParagraphNoStar}[1]{\oldsubparagraph{#1}\mbox{}}
\fi
\makeatother

\usepackage{color}
\usepackage{fancyvrb}
\newcommand{\VerbBar}{|}
\newcommand{\VERB}{\Verb[commandchars=\\\{\}]}
\DefineVerbatimEnvironment{Highlighting}{Verbatim}{commandchars=\\\{\}}
% Add ',fontsize=\small' for more characters per line
\usepackage{framed}
\definecolor{shadecolor}{RGB}{241,243,245}
\newenvironment{Shaded}{\begin{snugshade}}{\end{snugshade}}
\newcommand{\AlertTok}[1]{\textcolor[rgb]{0.68,0.00,0.00}{#1}}
\newcommand{\AnnotationTok}[1]{\textcolor[rgb]{0.37,0.37,0.37}{#1}}
\newcommand{\AttributeTok}[1]{\textcolor[rgb]{0.40,0.45,0.13}{#1}}
\newcommand{\BaseNTok}[1]{\textcolor[rgb]{0.68,0.00,0.00}{#1}}
\newcommand{\BuiltInTok}[1]{\textcolor[rgb]{0.00,0.23,0.31}{#1}}
\newcommand{\CharTok}[1]{\textcolor[rgb]{0.13,0.47,0.30}{#1}}
\newcommand{\CommentTok}[1]{\textcolor[rgb]{0.37,0.37,0.37}{#1}}
\newcommand{\CommentVarTok}[1]{\textcolor[rgb]{0.37,0.37,0.37}{\textit{#1}}}
\newcommand{\ConstantTok}[1]{\textcolor[rgb]{0.56,0.35,0.01}{#1}}
\newcommand{\ControlFlowTok}[1]{\textcolor[rgb]{0.00,0.23,0.31}{\textbf{#1}}}
\newcommand{\DataTypeTok}[1]{\textcolor[rgb]{0.68,0.00,0.00}{#1}}
\newcommand{\DecValTok}[1]{\textcolor[rgb]{0.68,0.00,0.00}{#1}}
\newcommand{\DocumentationTok}[1]{\textcolor[rgb]{0.37,0.37,0.37}{\textit{#1}}}
\newcommand{\ErrorTok}[1]{\textcolor[rgb]{0.68,0.00,0.00}{#1}}
\newcommand{\ExtensionTok}[1]{\textcolor[rgb]{0.00,0.23,0.31}{#1}}
\newcommand{\FloatTok}[1]{\textcolor[rgb]{0.68,0.00,0.00}{#1}}
\newcommand{\FunctionTok}[1]{\textcolor[rgb]{0.28,0.35,0.67}{#1}}
\newcommand{\ImportTok}[1]{\textcolor[rgb]{0.00,0.46,0.62}{#1}}
\newcommand{\InformationTok}[1]{\textcolor[rgb]{0.37,0.37,0.37}{#1}}
\newcommand{\KeywordTok}[1]{\textcolor[rgb]{0.00,0.23,0.31}{\textbf{#1}}}
\newcommand{\NormalTok}[1]{\textcolor[rgb]{0.00,0.23,0.31}{#1}}
\newcommand{\OperatorTok}[1]{\textcolor[rgb]{0.37,0.37,0.37}{#1}}
\newcommand{\OtherTok}[1]{\textcolor[rgb]{0.00,0.23,0.31}{#1}}
\newcommand{\PreprocessorTok}[1]{\textcolor[rgb]{0.68,0.00,0.00}{#1}}
\newcommand{\RegionMarkerTok}[1]{\textcolor[rgb]{0.00,0.23,0.31}{#1}}
\newcommand{\SpecialCharTok}[1]{\textcolor[rgb]{0.37,0.37,0.37}{#1}}
\newcommand{\SpecialStringTok}[1]{\textcolor[rgb]{0.13,0.47,0.30}{#1}}
\newcommand{\StringTok}[1]{\textcolor[rgb]{0.13,0.47,0.30}{#1}}
\newcommand{\VariableTok}[1]{\textcolor[rgb]{0.07,0.07,0.07}{#1}}
\newcommand{\VerbatimStringTok}[1]{\textcolor[rgb]{0.13,0.47,0.30}{#1}}
\newcommand{\WarningTok}[1]{\textcolor[rgb]{0.37,0.37,0.37}{\textit{#1}}}

\usepackage{longtable,booktabs,array}
\newcounter{none} % for unnumbered tables
\usepackage{calc} % for calculating minipage widths
% Correct order of tables after \paragraph or \subparagraph
\usepackage{etoolbox}
\makeatletter
\patchcmd\longtable{\par}{\if@noskipsec\mbox{}\fi\par}{}{}
\makeatother
% Allow footnotes in longtable head/foot
\IfFileExists{footnotehyper.sty}{\usepackage{footnotehyper}}{\usepackage{footnote}}
\makesavenoteenv{longtable}
\usepackage{graphicx}
\makeatletter
\newsavebox\pandoc@box
\newcommand*\pandocbounded[1]{% scales image to fit in text height/width
  \sbox\pandoc@box{#1}%
  \Gscale@div\@tempa{\textheight}{\dimexpr\ht\pandoc@box+\dp\pandoc@box\relax}%
  \Gscale@div\@tempb{\linewidth}{\wd\pandoc@box}%
  \ifdim\@tempb\p@<\@tempa\p@\let\@tempa\@tempb\fi% select the smaller of both
  \ifdim\@tempa\p@<\p@\scalebox{\@tempa}{\usebox\pandoc@box}%
  \else\usebox{\pandoc@box}%
  \fi%
}
% Set default figure placement to htbp
\def\fps@figure{htbp}
\makeatother





\setlength{\emergencystretch}{3em} % prevent overfull lines

\providecommand{\tightlist}{%
  \setlength{\itemsep}{0pt}\setlength{\parskip}{0pt}}



 


\usepackage{siunitx}
\usepackage{booktabs}
\usepackage{forest}
\usepackage{longtable}
\usepackage{xcolor}
\usepackage{fancyhdr}
\usepackage{titlesec}
\usepackage{caption}
\usepackage{subcaption}
\usepackage{enumitem}
\usepackage{graphicx}
\usepackage{hyperref}
\usepackage{listings}
\usepackage{array}
\usepackage{multirow}
% ── Color palette ──
\definecolor{accentblue}{HTML}{2D5F8A}
\definecolor{softgray}{HTML}{666666}
% ── Page style ──
\pagestyle{fancy}
\fancyhf{}
\fancyhead[C]{\small Ω-Architect 用户指南}
\fancyhead[R]{\small\leftmark}
\fancyfoot[C]{\thepage}
\renewcommand{\headrulewidth}{0.4pt}
\setlength{\headheight}{15pt}
% ── Section formatting ──
\titleformat{\chapter}{\Large\bfseries\color{accentblue}}{第\,\thechapter\,章}{1em}{}
\titleformat{\section}{\large\bfseries}{\S\thesection}{1em}{}
\titleformat{\subsection}{\normalsize\bfseries}{\S\thesubsection}{1em}{}
% ── Table styling ──
\captionsetup[table]{position=above,skip=3pt}
\captionsetup[figure]{position=below,skip=3pt}
\renewcommand{\arraystretch}{1.15}
% ── Code listing ──
\lstset{basicstyle=\small\ttfamily,breaklines=true,
  frame=single,numbers=left,numberstyle=\tiny\color{softgray}}
% ── Version info (generated by metadata script) ──
\IfFileExists{version-info.tex}{\input{version-info}}{}
\makeatletter
\@ifpackageloaded{bookmark}{}{\usepackage{bookmark}}
\makeatother
\makeatletter
\@ifpackageloaded{caption}{}{\usepackage{caption}}
\AtBeginDocument{%
\ifdefined\contentsname
  \renewcommand*\contentsname{Table of contents}
\else
  \newcommand\contentsname{Table of contents}
\fi
\ifdefined\listfigurename
  \renewcommand*\listfigurename{List of Figures}
\else
  \newcommand\listfigurename{List of Figures}
\fi
\ifdefined\listtablename
  \renewcommand*\listtablename{List of Tables}
\else
  \newcommand\listtablename{List of Tables}
\fi
\ifdefined\figurename
  \renewcommand*\figurename{Figure}
\else
  \newcommand\figurename{Figure}
\fi
\ifdefined\tablename
  \renewcommand*\tablename{Table}
\else
  \newcommand\tablename{Table}
\fi
}
\@ifpackageloaded{float}{}{\usepackage{float}}
\floatstyle{ruled}
\@ifundefined{c@chapter}{\newfloat{codelisting}{h}{lop}}{\newfloat{codelisting}{h}{lop}[chapter]}
\floatname{codelisting}{Listing}
\newcommand*\listoflistings{\listof{codelisting}{List of Listings}}
\makeatother
\makeatletter
\makeatother
\makeatletter
\@ifpackageloaded{caption}{}{\usepackage{caption}}
\@ifpackageloaded{subcaption}{}{\usepackage{subcaption}}
\makeatother
\usepackage{bookmark}
\IfFileExists{xurl.sty}{\usepackage{xurl}}{} % add URL line breaks if available
\urlstyle{same}
\hypersetup{
  pdftitle={Ω-Architect 用户指南},
  pdfauthor={Ω-Architect Team},
  pdfkeywords={Omega, Lean4, 形式化证明, 用户指南},
  colorlinks=true,
  linkcolor={blue},
  filecolor={Maroon},
  citecolor={blue!70!black},
  urlcolor={teal},
  pdfcreator={LaTeX via pandoc}}


\title{Ω-Architect 用户指南}
\usepackage{etoolbox}
\makeatletter
\providecommand{\subtitle}[1]{% add subtitle to \maketitle
  \apptocmd{\@title}{\par {\large #1 \par}}{}{}
}
\makeatother
\subtitle{面向 AI4Math 的端到端形式化证明系统 --- 安装、配置与接入}
\author{Ω-Architect Team}
\date{2026-06-10}
\begin{document}
\maketitle

\renewcommand*\contentsname{目录}
{
\setcounter{tocdepth}{2}
\tableofcontents
}

\bookmarksetup{startatroot}

\chapter{Ω-Architect
用户指南}\label{ux3c9-architect-ux7528ux6237ux6307ux5357}

\section{面向 AI4Math
的端到端形式化证明系统}\label{ux9762ux5411-ai4math-ux7684ux7aefux5230ux7aefux5f62ux5f0fux5316ux8bc1ux660eux7cfbux7edf}

\vspace{0.5cm}

\textbf{摘要：} Ω-Architect 是一个基于 Lean 4 + Mathlib
的开源形式化证明系统，面向 AI4Math
算法组、数学物理方程研究组、量子光学实验物理组。本文档提供系统架构概览、性能指标、快速入门指南、用户接入说明、国内外研究对比，以及物理/量子领域到形式化数学的交叉前景。

\vspace{0.3cm}

{\def\LTcaptype{none} % do not increment counter
\begin{longtable}[]{@{}ll@{}}
\toprule\noalign{}
\endhead
\bottomrule\noalign{}
\endlastfoot
\textbf{文档版本} & \texttt{v0.1.0} (commit
\texttt{\textbackslash{}texttt\{2953191\}}) \\
\textbf{提交日期} & \texttt{2026-06-10} \\
\textbf{repo 版本} & \texttt{v0.1.0}（与 \texttt{pyproject.toml}
同步） \\
\textbf{许可证} & Apache 2.0 \\
\textbf{目标} & MiniF2F 244 题通过率 \textgreater95\% \\
\textbf{当前开发集} & Easy 3/3 ✅, Medium 1/4, Hard 0/3 \\
\end{longtable}
}

\textbf{阅读指引：}

{\def\LTcaptype{none} % do not increment counter
\begin{longtable}[]{@{}ll@{}}
\toprule\noalign{}
读者群体 & 推荐章节 \\
\midrule\noalign{}
\endhead
\bottomrule\noalign{}
\endlastfoot
AI4Math 算法组 & ch04 Quick Start, ch02 系统架构 \\
数学物理方程研究组 & ch05 交叉领域, ch06 路线图 \\
量子光学实验物理组 & ch05 量子光学需求, appx 验证实例 \\
\end{longtable}
}

\newpage

\bookmarksetup{startatroot}

\chapter{项目概述与定位}\label{ux9879ux76eeux6982ux8ff0ux4e0eux5b9aux4f4d}

\section{技术定位}\label{sec-positioning}

Ω-Architect 处于 AI 形式化数学验证的技术交叉点：

\begin{verbatim}
大语言模型(LLM) ⟶ 证明策略生成 ⟶ Lean 4 编译器验证
\end{verbatim}

它与以下方向有本质区别：

{\def\LTcaptype{none} % do not increment counter
\begin{longtable}[]{@{}lll@{}}
\toprule\noalign{}
方向 & 代表工作 & 区别 \\
\midrule\noalign{}
\endhead
\bottomrule\noalign{}
\endlastfoot
纯 LLM 推理 & GPT-4o, DeepSeek & 无形式化验证，存在幻觉 \\
传统自动化定理证明 & Eprover, Vampire & 不支持高阶逻辑、Mathlib \\
交互式证明助手 & Lean 4, Coq, Isabelle & 需人工编写所有证明 \\
RL-based 证明 & AlphaProof (DeepMind) & 闭源、专用算力 \\
\end{longtable}
}

\section{核心设计原则}\label{sec-design-principles}

\begin{enumerate}
\def\labelenumi{\arabic{enumi}.}
\tightlist
\item
  \textbf{开源优先}：Apache 2.0 许可证，全链路开放
\item
  \textbf{T2 验证为真值来源}：只有通过真实 Lean 编译器验证的证明才算有效
\item
  \textbf{多引擎容错}：Goedel + Rethlas + Archon
  三路并行，任一失败自动切换
\item
  \textbf{ML 驱动的自适应路由}：LightGBM
  模型自动判断定理难度，分配最优证明引擎
\item
  \textbf{预算感知}：BudgetTracker 自动管理
  token/时间/尝试次数，防止无限消耗
\end{enumerate}

\section{与现有系统的关系}\label{sec-relationship}

\begin{forest}
for tree={grow=east,edge={gray!60},l sep=6pt,s sep=8pt,
  font=\footnotesize,grow'=east}
[Ω-Architect
  [{Goedel-Prover / 并行+自修正}]
  [{Rethlas / 蓝图分解}]
  [{Archon / 迭代精炼}]
  [{DeepSeek API}]
  [{vLLM / 本地推理}]
  [{Mathlib / 170万定理}]
]
\end{forest}

\newpage

\bookmarksetup{startatroot}

\chapter{系统架构与核心模块}\label{ux7cfbux7edfux67b6ux6784ux4e0eux6838ux5fc3ux6a21ux5757}

\section{目录结构}\label{sec-directory}

\begin{verbatim}
omega-architect/
├── omega/                      # 核心源码
│   ├── __init__.py             # 版本号 (0.1.0)
│   ├── llm.py                  # LLM 连接层 (DeepSeek / vLLM / Ollama)
│   ├── runner.py               # OmegaRunner — 主入口
│   ├── data.py                 # 数据定义与 JSONL 实验记录
│   ├── prover/                 # 证明器三件套
│   │   ├── go_prover.py        # GoedelProver — 并行采样+自修正
│   │   ├── re_prover.py        # RethlasProver — 蓝图分解
│   │   ├── ar_prover.py        # ArchonProver — 迭代精炼
│   │   ├── ensemble.py         # EnsembleProver — 三路集成
│   │   ├── playbook.py         # ACE 风格策略积累
│   │   └── cache.py            # 证明缓存
│   ├── search/                 # 搜索与提议
│   │   ├── proposer.py         # Proposer — 策略生成 (LLM + 模板)
│   │   ├── tree.py             # GoalState / SearchNode
│   │   ├── blueprint.py        # 蓝图分解
│   │   ├── error_classifier.py # Lean 编译错误 13 类分类
│   │   ├── curriculum.py       # 课程调度
│   │   ├── lean_search.py      # Lean 搜索集成
│   │   ├── matlas_cache.py     # Mathlib 缓存管理
│   │   └── passk.py            # Pass@K 统计
│   ├── verify/                 # 验证层
│   │   ├── t1_llm.py           # T1: LLM 快速验证 (~5s)
│   │   ├── t2_lean.py          # T2: 诊断解析
│   │   ├── t2_real.py          # T2: lake env lean --stdin 真实编译
│   │   └── t2_mcp.py           # T2: MCP 编译通道
│   ├── resource/               # 资源与路由
│   │   ├── model_router.py     # ModelRouter + MLRoute
│   │   ├── routing_flags.py    # 路由标志系统
│   │   ├── budget.py           # BudgetTracker
│   │   ├── config.py           # BudgetConfig
│   │   ├── pricing.py          # 模型定价表
│   │   ├── allocator.py        # 模型分配器
│   │   └── models/             # LightGBM/ONNX 模型文件
│   ├── loop/                   # 循环工程
│   │   ├── inner.py            # Inner Loop — tool_calls + thinking
│   │   ├── runner.py           # LoopExperiment
│   │   ├── gate.py             # ConvergenceTracker
│   │   ├── errors.py           # 编译错误分类
│   │   └── dialogue_cache.py   # 成功对话缓存
│   ├── gpu_layer/              # GPU 后端自动路由
│   │   ├── detector.py         # 硬件检测
│   │   ├── backends.py         # vLLM / Ollama / Transformers 统一接口
│   │   └── scheduler.py        # 全局调度器
│   ├── benchmark/              # 基准测试
│   │   └── suite.py            # 4 级难度金字塔
│   ├── cli/                    # CLI 命令
│   ├── agent/                  # Agent 编排层
│   │   ├── orchestrator.py     # 状态机路由
│   │   └── message_log.py      # 消息日志
│   └── search/                 # 搜索与提议
├── tests/                      # 测试套件 (25 文件, 549 测试)
├── docs/                       # 文档
│   ├── whitepaper/             # 技术白皮书 (Quarto)
│   ├── user_guide/             # 用户指南 (本文档)
│   └── design/                 # 设计文档
├── scripts/                    # 运维脚本
├── benchmarks/minif2f/         # MiniF2F 基准
└── pyproject.toml              # 项目配置 (版本 0.1.0)
\end{verbatim}

\section{核心数据流}\label{sec-dataflow}

\begin{verbatim}
[定理输入]
    │
    ▼
ModelRouter (analyze_theorem_pattern)
    │
    ▼
MLRoute (LightGBM 95.8% 准确率, <100µs)
    │
    ▼
EnsembleProver
├── GoedelProver  → vLLM (Goedel-Prover-V2-8B 本地)
├── RethlasProver → LLM (蓝图分解)
└── ArchonProver  → LLM (迭代精炼)
    │
    ▼
T1 Verify (LLM, ~5s)
    │
    ▼
T2 Real Compile (lake env lean --stdin, ~2.5s)
    │
    ▼
[已验证的 Lean 证明 → GoedelResult]
\end{verbatim}

\section{关键接口}\label{sec-interfaces}

{\def\LTcaptype{none} % do not increment counter
\begin{longtable}[]{@{}lll@{}}
\toprule\noalign{}
接口 & 位置 & 功能 \\
\midrule\noalign{}
\endhead
\bottomrule\noalign{}
\endlastfoot
\texttt{OmegaRunner.prove()} & \texttt{runner.py} & 单定理证明入口 \\
\texttt{resolve\_generate\_fn()} & \texttt{llm.py} & 模型路由选择 \\
\texttt{GoedelProver.run()} & \texttt{prover/go\_prover.py} & Goedel
式并行证明 \\
\texttt{EnsembleProver.run()} & \texttt{prover/ensemble.py} &
三路集成证明 \\
\texttt{make\_real\_compile\_callback()} & \texttt{verify/t2\_real.py} &
T2 编译回调 \\
\texttt{ModelRouter.select()} & \texttt{resource/model\_router.py} &
模型智能选择 \\
\texttt{BudgetTracker} & \texttt{resource/budget.py} & 预算跟踪 \\
\end{longtable}
}

\section{验证层架构}\label{sec-verification}

\begin{verbatim}
T1 (LLM) ── 快速检查, ~5s, 捕获明显错误
   │
   ▼ (通过)
T2 (Real Compile) ── lake env lean --stdin, ~2.5s
   │               └── Mathlib 8109 oleans, 6.7GB
   ▼ (exit_code=0)
[已验证通过]
\end{verbatim}

\textbf{T2 失败时的回退链：}

\begin{enumerate}
\def\labelenumi{\arabic{enumi}.}
\tightlist
\item
  错误解析 → 13 类分类
\item
  自适应策略注入 (Auto aesop / simp / nlinarith / ring / positivity)
\item
  子引理自动生长（unknown identifier 时展开新节点）
\item
  MCP \texttt{lean\_multi\_attempt} 多策略并行尝试
\item
  对话保存到 \texttt{dialogue\_cache.py} 供后续学习
\end{enumerate}

\newpage

\bookmarksetup{startatroot}

\chapter{技术性能指标}\label{ux6280ux672fux6027ux80fdux6307ux6807}

\section{测试覆盖}\label{sec-test-coverage}

{\def\LTcaptype{none} % do not increment counter
\begin{longtable}[]{@{}
  >{\raggedright\arraybackslash}p{(\linewidth - 2\tabcolsep) * \real{0.5000}}
  >{\raggedright\arraybackslash}p{(\linewidth - 2\tabcolsep) * \real{0.5000}}@{}}
\toprule\noalign{}
\begin{minipage}[b]{\linewidth}\raggedright
指标
\end{minipage} & \begin{minipage}[b]{\linewidth}\raggedright
数值
\end{minipage} \\
\midrule\noalign{}
\endhead
\bottomrule\noalign{}
\endlastfoot
测试文件数 & 25 \\
总测试数 & 549 \\
通过率 & 99.8\% (1 个预存失败，无关) \\
覆盖模块 & prover / verify / search / resource / llm / loop /
gpu\_layer \\
\end{longtable}
}

\section{编译性能}\label{sec-compile-performance}

{\def\LTcaptype{none} % do not increment counter
\begin{longtable}[]{@{}lll@{}}
\toprule\noalign{}
操作 & 耗时 & 备注 \\
\midrule\noalign{}
\endhead
\bottomrule\noalign{}
\endlastfoot
T2 编译（含 Mathlib） & \textasciitilde2.5 s/定理 &
\texttt{lake\ env\ lean\ -\/-stdin} \\
T2 编译（无 Mathlib） & \textasciitilde85 ms & Init-only \\
T1 LLM 验证 & \textasciitilde5 s & 单次 LLM 调用 \\
JSON mode 推理 (DeepSeek) & \textasciitilde2 s &
\texttt{response\_format=json\_object} \\
vLLM 推理 (Goedel-8B) & \textasciitilde8 s & 本地 24.38 tok/s \\
MLRoute 推理 & \textless100 µs & LightGBM + ONNX \\
\end{longtable}
}

\section{模型定价与路由}\label{sec-pricing}

{\def\LTcaptype{none} % do not increment counter
\begin{longtable}[]{@{}llll@{}}
\toprule\noalign{}
模型 & 输入 \$/M tok & 输出 \$/M tok & 适用难度 \\
\midrule\noalign{}
\endhead
\bottomrule\noalign{}
\endlastfoot
Goedel-Prover-V2-8B (本地) & 0 & 0 & Simple / Medium \\
DeepSeek V4 Flash & 0.14 & 0.28 & Medium / Hard \\
DeepSeek V4 Pro & 0.42 & 0.84 & Hard \\
Qwen3-Coder:30b (Ollama) & 0 & 0 & Simple \\
\end{longtable}
}

对于典型定理（\textasciitilde1000 token 输入 / \textasciitilde500 token
输出），DeepSeek Flash 调用的成本约 \textbf{\$0.0005}。

\section{MiniF2F 基准进展}\label{sec-benchmark}

{\def\LTcaptype{none} % do not increment counter
\begin{longtable}[]{@{}
  >{\raggedright\arraybackslash}p{(\linewidth - 6\tabcolsep) * \real{0.2000}}
  >{\raggedright\arraybackslash}p{(\linewidth - 6\tabcolsep) * \real{0.3000}}
  >{\raggedright\arraybackslash}p{(\linewidth - 6\tabcolsep) * \real{0.3000}}
  >{\raggedright\arraybackslash}p{(\linewidth - 6\tabcolsep) * \real{0.2000}}@{}}
\toprule\noalign{}
\begin{minipage}[b]{\linewidth}\raggedright
基准
\end{minipage} & \begin{minipage}[b]{\linewidth}\raggedright
T1 通过
\end{minipage} & \begin{minipage}[b]{\linewidth}\raggedright
T2 通过
\end{minipage} & \begin{minipage}[b]{\linewidth}\raggedright
状态
\end{minipage} \\
\midrule\noalign{}
\endhead
\bottomrule\noalign{}
\endlastfoot
MiniF2F (50 sample) & 50/50 (100\%) & 0/50 & 需 Goedel-Prover
全量运行 \\
内部 Tier 1 (纯 Lean) & 5/5 & --- & 基础设施就绪 \\
内部 Tier 2 (Mathlib simp) & 10/10 & --- & 基础设施就绪 \\
内部 Tier 3 (Induction) & 10/10 & --- & 基础设施就绪 \\
内部 Tier 4 (MiniF2F 子集) & 5/5 & --- & 基础设施就绪 \\
\end{longtable}
}

\textbf{注：} T2 验证需 Goedel-Prover-V2-8B vLLM 服务运行中（本地 GPU
24GB）。当前 T1 验证已完全就绪（100\% 通过），T2
单例验证已通过（\texttt{1+1=2}、\texttt{mul\_zero}）。

\newpage

\bookmarksetup{startatroot}

\chapter{用户 Quick Start
与进阶管线}\label{ux7528ux6237-quick-start-ux4e0eux8fdbux9636ux7ba1ux7ebf}

\section{快速开始（5 分钟）}\label{sec-quickstart}

\begin{Shaded}
\begin{Highlighting}[]
\CommentTok{\# 1. 克隆}
\FunctionTok{git}\NormalTok{ clone https://github.com/your{-}org/omega{-}architect.git}
\BuiltInTok{cd}\NormalTok{ omega{-}architect}

\CommentTok{\# 2. 安装}
\ExtensionTok{pip}\NormalTok{ install }\AttributeTok{{-}e}\NormalTok{ .}
\ExtensionTok{omega}\NormalTok{ init }\AttributeTok{{-}{-}force}   \CommentTok{\# 检测 Lean / Mathlib}

\CommentTok{\# 3. 证明一个定理}
\ExtensionTok{omega}\NormalTok{ prove }\StringTok{"theorem t : 1 + 1 = 2 :="} \AttributeTok{{-}{-}model}\NormalTok{ deepseek/deepseek{-}v4{-}flash}
\end{Highlighting}
\end{Shaded}

输出：

\begin{verbatim}
[OK] proof found in 8.23s (1 attempts, 0 correction rounds)
\end{verbatim}

\section{环境要求}\label{sec-requirements}

{\def\LTcaptype{none} % do not increment counter
\begin{longtable}[]{@{}lll@{}}
\toprule\noalign{}
组件 & 最低要求 & 推荐 \\
\midrule\noalign{}
\endhead
\bottomrule\noalign{}
\endlastfoot
Python & 3.11 & 3.13 \\
GPU & --- & RTX 4500 Ada 24GB \\
RAM & 16 GB & 32 GB \\
磁盘 & 10 GB & 20 GB (含 Mathlib 缓存 6.7 GB) \\
Lean 4 & 4.30.0 & 4.30.0+ \\
Mathlib & 8109 oleans & 完整缓存 \\
\end{longtable}
}

\section{进阶开发管线}\label{sec-advanced}

\subsection{使用本地
Goedel-Prover-V2}\label{ux4f7fux7528ux672cux5730-goedel-prover-v2}

\begin{Shaded}
\begin{Highlighting}[]
\CommentTok{\# 启动 vLLM 服务器}
\FunctionTok{bash}\NormalTok{ scripts/start{-}goedel.sh}

\CommentTok{\# 使用本地模型}
\ExtensionTok{omega}\NormalTok{ prove }\StringTok{"..."} \AttributeTok{{-}{-}model}\NormalTok{ goedel/goedel{-}v2{-}8b}
\end{Highlighting}
\end{Shaded}

\subsection{运行基准测试}\label{ux8fd0ux884cux57faux51c6ux6d4bux8bd5}

\begin{Shaded}
\begin{Highlighting}[]
\CommentTok{\# 按难度层级运行}
\ExtensionTok{omega}\NormalTok{ benchmark }\AttributeTok{{-}{-}tier}\NormalTok{ 1{-}3 }\AttributeTok{{-}{-}model}\NormalTok{ deepseek/deepseek{-}v4{-}flash}

\CommentTok{\# 自定义预算}
\ExtensionTok{omega}\NormalTok{ prove }\StringTok{"..."} \AttributeTok{{-}{-}budget{-}time}\NormalTok{ 600 }\AttributeTok{{-}{-}budget{-}cost}\NormalTok{ 0.5}

\CommentTok{\# 查看配置}
\ExtensionTok{omega}\NormalTok{ config show }\AttributeTok{{-}{-}lean}
\end{Highlighting}
\end{Shaded}

\section{编程接口}\label{sec-api}

\begin{Shaded}
\begin{Highlighting}[]
\ImportTok{from}\NormalTok{ omega.prover }\ImportTok{import}\NormalTok{ GoedelProver}
\ImportTok{from}\NormalTok{ omega.verify.t2\_real }\ImportTok{import}\NormalTok{ make\_real\_compile\_callback}

\CommentTok{\# 初始化编译器}
\NormalTok{compile\_fn }\OperatorTok{=}\NormalTok{ make\_real\_compile\_callback()}

\CommentTok{\# 创建证明器}
\NormalTok{prover }\OperatorTok{=}\NormalTok{ GoedelProver(compile\_fn}\OperatorTok{=}\NormalTok{compile\_fn, num\_samples}\OperatorTok{=}\DecValTok{6}\NormalTok{)}

\CommentTok{\# 运行证明}
\NormalTok{result }\OperatorTok{=}\NormalTok{ prover.run(}
    \StringTok{"theorem mul\_zero (a : Nat) : a * 0 = 0 :="}
\NormalTok{)}
\BuiltInTok{print}\NormalTok{(result.summary)  }\CommentTok{\# "[OK] proof found in 12.5s"}
\BuiltInTok{print}\NormalTok{(result.proof)    }\CommentTok{\# 编译通过的完整 Lean 代码}
\end{Highlighting}
\end{Shaded}

\section{OmegaRunner 长时间运行}\label{sec-runner}

\begin{Shaded}
\begin{Highlighting}[]
\ImportTok{from}\NormalTok{ omega.runner }\ImportTok{import}\NormalTok{ OmegaRunner}

\NormalTok{runner }\OperatorTok{=}\NormalTok{ OmegaRunner(time\_budget\_s}\OperatorTok{=}\DecValTok{14400}\NormalTok{)  }\CommentTok{\# 4 小时}
\NormalTok{report }\OperatorTok{=}\NormalTok{ runner.run(}
    \StringTok{"Formalize NV center spin Hamiltonian"}
\NormalTok{)}
\BuiltInTok{print}\NormalTok{(report.summary())}
\end{Highlighting}
\end{Shaded}

支持 checkpoint/resume，适合 2/4/8/12 小时无人值守证明搜索会话。

\newpage

\bookmarksetup{startatroot}

\chapter{数据集、国内外对比与交叉领域}\label{ux6570ux636eux96c6ux56fdux5185ux5916ux5bf9ux6bd4ux4e0eux4ea4ux53c9ux9886ux57df}

\section{AI4Math 证明题数据集对比}\label{sec-datasets}

{\def\LTcaptype{none} % do not increment counter
\begin{longtable}[]{@{}lllll@{}}
\toprule\noalign{}
数据集 & 题数 & 难度等级 & 支撑系统 & Lean 4 支持 \\
\midrule\noalign{}
\endhead
\bottomrule\noalign{}
\endlastfoot
MiniF2F & 488 & IMO \textasciitilde{} AMC & Lean / Isabelle / Metamath &
✅ 完整 \\
PutnamBench & 640 & Putnam 竞赛 & Lean 4 & ✅ \\
ProofNet & 371 & 本科数学 & Lean 3 / 4 & 🔄 迁移中 \\
Lean-workbook & 59,880 & 合成+数学竞赛 & Lean 4 & ✅ \\
SorryDB & 1,000+ & 真实世界开源项目 & Lean 4 & ✅ \\
FormalNumina & 120,000+ & 合成（基于 Numina） & Lean 4 & ✅ \\
IMO Shortlist & \textasciitilde300 & IMO 级 & --- & ⬜ 部分 \\
FATE (PKU) & 200 & 博士级代数 & Lean 4 & ✅ \\
\end{longtable}
}

\section{国内外相关研究工作}\label{sec-related-work}

\subsection{国际}\label{ux56fdux9645}

{\def\LTcaptype{none} % do not increment counter
\begin{longtable}[]{@{}
  >{\raggedright\arraybackslash}p{(\linewidth - 6\tabcolsep) * \real{0.2143}}
  >{\raggedright\arraybackslash}p{(\linewidth - 6\tabcolsep) * \real{0.2143}}
  >{\raggedright\arraybackslash}p{(\linewidth - 6\tabcolsep) * \real{0.2143}}
  >{\raggedright\arraybackslash}p{(\linewidth - 6\tabcolsep) * \real{0.3571}}@{}}
\toprule\noalign{}
\begin{minipage}[b]{\linewidth}\raggedright
工作
\end{minipage} & \begin{minipage}[b]{\linewidth}\raggedright
机构
\end{minipage} & \begin{minipage}[b]{\linewidth}\raggedright
年份
\end{minipage} & \begin{minipage}[b]{\linewidth}\raggedright
核心方法
\end{minipage} \\
\midrule\noalign{}
\endhead
\bottomrule\noalign{}
\endlastfoot
AlphaProof & DeepMind & 2024-25 & RL + 8000 万形式化问题，IMO 2025
金牌 \\
Goedel-Prover-V2 & Goedel-LM & 2025 & 脚手架数据合成+自修正，MiniF2F
88.1\% \\
DeepSeek-Prover-V2 & DeepSeek & 2025 & 671B MoE，RL 训练，PutnamBench 47
题 \\
LeanDojo-v2 & Caltech & 2025 & 检索增强 LLM + 仓库抽取框架 \\
LeanAgent & Caltech & 2025 & 终身学习+迁移证明策略 \\
SorryDB & 多机构 & 2026 & 实时开源 Lean 证明任务基准 \\
Archon & 多机构 & 2026 & 自动猜想解析+形式化验证 \\
FormalProofBench & --- & 2026 & 研究生级数学证明基准 \\
Mathesis & 华为 & 2025 & NL→Lean4 自动形式化 \\
\end{longtable}
}

\subsection{国内}\label{ux56fdux5185}

{\def\LTcaptype{none} % do not increment counter
\begin{longtable}[]{@{}lll@{}}
\toprule\noalign{}
工作 & 机构 & 方向 \\
\midrule\noalign{}
\endhead
\bottomrule\noalign{}
\endlastfoot
FATE (200 题) & 北京大学 & 博士级代数定理数据集 \\
Herald (580K NL↔FL) & PKU FrenzyMath & 自然语言-形式语言配对 \\
OProofs (6.86M 证明) & --- & 大量形式化证明语料 \\
Rethlas & PKU FrenzyMath & 蓝图分解+回溯证明搜索 \\
DeepSeek-Prover 系列 & 深度求索 & 671B MoE 证明模型 \\
\end{longtable}
}

\subsection{性能对比（MiniF2F
Pass@32）}\label{ux6027ux80fdux5bf9ux6bd4minif2f-pass32}

{\def\LTcaptype{none} % do not increment counter
\begin{longtable}[]{@{}ll@{}}
\toprule\noalign{}
系统 & 通过率 \\
\midrule\noalign{}
\endhead
\bottomrule\noalign{}
\endlastfoot
Goedel-Prover-V2-32B & 88.1\% \\
Goedel-Prover-V2-8B & 84.6\% \\
DeepSeek-Prover-V2 & \textasciitilde80\% \\
Goedel-Prover-V1 & 57.6\% \\
\textbf{Omega（当前）} & \textbf{T1 100\%, T2 E2E 验证中} \\
\end{longtable}
}

\section{交叉领域：物理、量子光学到形式化数学}\label{sec-cross-domain}

\subsection{物理数学形式化的现状}\label{ux7269ux7406ux6570ux5b66ux5f62ux5f0fux5316ux7684ux73b0ux72b6}

目前 Lean 4 / Mathlib
主要集中在纯数学领域。物理领域的形式化进展相对早期：

{\def\LTcaptype{none} % do not increment counter
\begin{longtable}[]{@{}lll@{}}
\toprule\noalign{}
领域 & 已有工作 & 成熟度 \\
\midrule\noalign{}
\endhead
\bottomrule\noalign{}
\endlastfoot
泛函分析 & Mathlib 4 & ✅ 完善 \\
张量代数 & Mathlib 4 & ✅ 完善 \\
微分几何 & Mathlib 4 & ✅ 完善 \\
微分方程 & 有限（ODE 基本） & 🔄 起步 \\
电磁学 & 极少 & ⬜ \\
量子力学 & 极少 & ⬜ \\
量子光学 & 几乎为零 & ⬜ \\
神经网络形式化 & TorchLean & 🔄 进行中 \\
\end{longtable}
}

\subsection{SciLean 项目}\label{scilean-ux9879ux76ee}

\href{https://github.com/lecopivo/SciLean}{SciLean} 是目前最接近''物理用
Mathlib''的项目之一：

\begin{itemize}
\tightlist
\item
  在 Lean 4 中实现可微编程
\item
  支持向量微积分、张量运算
\item
  适合 PDE 的符号推导+数值验证
\item
  许可证：Apache 2.0
\end{itemize}

\subsection{Physlib：物理形式化的社区基座}\label{physlibux7269ux7406ux5f62ux5f0fux5316ux7684ux793eux533aux57faux5ea7}

\href{https://github.com/leanprover-community/physlib}{Physlib}（\href{https://physlib.io}{physlib.io}）是
Lean 4 社区官方物理形式化库，由原 PhysLean（通用物理）和
Lean-QuantumInfo（量子信息）合并而成：

\begin{itemize}
\tightlist
\item
  \textbf{目标}：成为物理学的
  Mathlib------统一、开源、社区驱动的物理定理库
\item
  \textbf{范围}：经典力学、电磁学、粒子物理、宇宙学、量子信息
\item
  \textbf{规模}：3,169 commits，Lean v4.30.0
\item
  \textbf{组织结构}：\texttt{Physlib/}（通用物理）+
  \texttt{QuantumInfo/}（量子信息），可分别构建
\item
  \textbf{社区}：Lean Zulip 专用频道（\#479953-Physlib），Gitpod
  即开即用
\item
  \textbf{论文}：arXiv:2405.08863（Joseph Tooby-Smith）
\item
  \textbf{许可证}：MIT / 开源
\end{itemize}

\textbf{QuantumInfo 子库}（原
\href{https://github.com/Timeroot/Lean-QuantumInfo}{Timeroot/Lean-QuantumInfo}，已合入
Physlib）：

\begin{itemize}
\tightlist
\item
  有限维 Hilbert 空间上的量子信息论形式化
\item
  涵盖 Mark Wilde《Quantum Information Theory》核心定义与定理
\item
  已完成的标志性定理：\textbf{Generalized Quantum Stein's
  Lemma}（\texttt{limit\_hypotesting\_eq\_limit\_rel\_entropy}）
\item
  配套研究计划：\href{https://qform.al}{QFormal}（Quantum Formal
  Initiative）------量子算法、量子密码学安全证明的形式化验证
\item
  论文：arXiv:2510.08672（Quantum Stein's Lemma 形式化）
\item
  衍生框架：arXiv:2510.26094（Lean4Physics/Lean4PHYS------基于 Physlib
  的大学物理推理框架）
\end{itemize}

\subsection{Ω-Architect
在交叉领域的定位}\label{ux3c9-architect-ux5728ux4ea4ux53c9ux9886ux57dfux7684ux5b9aux4f4d}

\begin{verbatim}
[物理问题] → [NL 描述] → [Mathesis / 自动形式化] → [Lean 4 定理]
                                                 ↓
[Physlib / SciLean 物理库] ← [Ω-Architect 自动证明] ← [Mathlib 数学基础]
\end{verbatim}

Omega 不直接构建物理形式化库，而是作为证明生成引擎插入这一链条：

\begin{itemize}
\tightlist
\item
  \textbf{输入}：Lean 4 格式的物理定理（基于 Physlib 或 SciLean 定义）
\item
  \textbf{输出}：T2 编译通过的证明
\item
  \textbf{物理库依赖}：Physlib（社区基座）+ SciLean（可微编程）+
  社区贡献
\end{itemize}

\textbf{Omega + Physlib 的协作模式：}

{\def\LTcaptype{none} % do not increment counter
\begin{longtable}[]{@{}
  >{\raggedright\arraybackslash}p{(\linewidth - 4\tabcolsep) * \real{0.3333}}
  >{\raggedright\arraybackslash}p{(\linewidth - 4\tabcolsep) * \real{0.3333}}
  >{\raggedright\arraybackslash}p{(\linewidth - 4\tabcolsep) * \real{0.3333}}@{}}
\toprule\noalign{}
\begin{minipage}[b]{\linewidth}\raggedright
角色
\end{minipage} & \begin{minipage}[b]{\linewidth}\raggedright
谁做
\end{minipage} & \begin{minipage}[b]{\linewidth}\raggedright
产出
\end{minipage} \\
\midrule\noalign{}
\endhead
\bottomrule\noalign{}
\endlastfoot
物理定义与公理 & Physlib 社区 & \([\hat{a}, \hat{a}^\dagger] = 1\)
的代数公理 \\
定理形式化 & 物理组 / Mathesis & 量子光学中的具体命题 \\
自动证明生成 & Ω-Architect & T2 编译通过的 Lean 证明 \\
验证与合并 & Physlib 审阅流程 & 回馈到 Physlib 主库 \\
\end{longtable}
}

\subsection{量子光学的特殊需求}\label{ux91cfux5b50ux5149ux5b66ux7684ux7279ux6b8aux9700ux6c42}

量子光学形式化需要：

\begin{enumerate}
\def\labelenumi{\arabic{enumi}.}
\tightlist
\item
  \textbf{算符代数}：\(\hat{a}, \hat{a}^\dagger\) 产生湮灭算符的交换关系
\item
  \textbf{希尔伯特空间}：态矢量的内积、张量积
\item
  \textbf{密度矩阵}：迹运算、偏迹
\item
  \textbf{主方程}：Lindblad 形式的微分方程
\end{enumerate}

这些在 Physlib 的 QuantumInfo 子库中已有基础（有限维 Hilbert
空间、密度算子、迹运算），但量子光学特有的连续变量体系、相干态、压缩态等仍需扩展。Omega
可以在已有量子信息公理基础上自动推导这些关系。

\subsection{已有开源工作}\label{ux5df2ux6709ux5f00ux6e90ux5de5ux4f5c}

{\def\LTcaptype{none} % do not increment counter
\begin{longtable}[]{@{}
  >{\raggedright\arraybackslash}p{(\linewidth - 4\tabcolsep) * \real{0.2500}}
  >{\raggedright\arraybackslash}p{(\linewidth - 4\tabcolsep) * \real{0.2500}}
  >{\raggedright\arraybackslash}p{(\linewidth - 4\tabcolsep) * \real{0.5000}}@{}}
\toprule\noalign{}
\begin{minipage}[b]{\linewidth}\raggedright
项目
\end{minipage} & \begin{minipage}[b]{\linewidth}\raggedright
领域
\end{minipage} & \begin{minipage}[b]{\linewidth}\raggedright
链接 / DOI
\end{minipage} \\
\midrule\noalign{}
\endhead
\bottomrule\noalign{}
\endlastfoot
Physlib & 物理形式化库（社区） & GitHub / physlib.io /
arXiv:2405.08863 \\
Lean-QuantumInfo / QFormal & 量子信息形式化 & GitHub / qform.al /
arXiv:2510.08672 \\
Lean4Physics & 大学物理推理框架 & arXiv:2510.26094 \\
SciLean & 科学计算/可微编程 & GitHub \\
TorchLean & 神经网络形式化 & arXiv:2602.22631 \\
Mathesis (华为) & NL→Lean4 形式化 & GitHub / arXiv:2506.07047 \\
LeanDojo-v2 & AI 辅助证明框架 & leandojo.org \\
Archon & 自动猜想解析 & arXiv:2604.03789 \\
Goedel-Prover-V2 & SOTA 证明模型 & GitHub / arXiv:2508.03613 \\
DeepSeek-Prover-V2 & 大型证明模型 & GitHub \\
\end{longtable}
}

\newpage

\bookmarksetup{startatroot}

\chapter{目标场景、开发状态与路线图}\label{ux76eeux6807ux573aux666fux5f00ux53d1ux72b6ux6001ux4e0eux8defux7ebfux56fe}

\section{目标场景清单}\label{sec-scenarios}

{\def\LTcaptype{none} % do not increment counter
\begin{longtable}[]{@{}
  >{\raggedright\arraybackslash}p{(\linewidth - 8\tabcolsep) * \real{0.1538}}
  >{\raggedright\arraybackslash}p{(\linewidth - 8\tabcolsep) * \real{0.1538}}
  >{\raggedright\arraybackslash}p{(\linewidth - 8\tabcolsep) * \real{0.2051}}
  >{\raggedright\arraybackslash}p{(\linewidth - 8\tabcolsep) * \real{0.1538}}
  >{\raggedright\arraybackslash}p{(\linewidth - 8\tabcolsep) * \real{0.3333}}@{}}
\toprule\noalign{}
\begin{minipage}[b]{\linewidth}\raggedright
场景
\end{minipage} & \begin{minipage}[b]{\linewidth}\raggedright
描述
\end{minipage} & \begin{minipage}[b]{\linewidth}\raggedright
优先级
\end{minipage} & \begin{minipage}[b]{\linewidth}\raggedright
状态
\end{minipage} & \begin{minipage}[b]{\linewidth}\raggedright
外部帮助需求
\end{minipage} \\
\midrule\noalign{}
\endhead
\bottomrule\noalign{}
\endlastfoot
数学竞赛证明 & IMO / AMC 级别问题 & P0 & 🔄 E2E 验证中 & 无 \\
本科数学定理 & 推导验证 & P0 & ✅ 管线已通 & 无 \\
论文公式验证 & 学术论文中定理的 Lean 形式化 & P1 & 📋 规划中 &
形式化专家 \\
物理方程推导 & 基于 Physlib 的 PDE/ODE 形式化 & P1 & 📋 Physlib
就绪，Omega 集成中 & 数学物理组 \\
量子光学计算 & 基于 QuantumInfo 的量子态演化、算符代数 & P2 & 📋
QuantumInfo 就绪，Omega 集成待启动 & 量子光学实验组 \\
自动形式化 & NL→Lean4 翻译 & P0 & 🔄 对接 Mathesis & 无 \\
ProofNet 覆盖 & 本科数学题库 & P1 & ✅ 数据就绪 & 无 \\
T2 大规模运行 & Goedel-Prover 批量证明 & P0 & ✅ GPU 资源充足 & GPU (RTX
4500 Ada) \\
\end{longtable}
}

\section{已就绪能力}\label{sec-ready}

\begin{itemize}
\tightlist
\item
  ✅ 端到端管线：LLM → T2 编译 → 验证通过
\item
  ✅ 3 路证明引擎集成（Goedel / Rethlas / Archon）
\item
  ✅ DeepSeek JSON mode（结构化输出+置信度）
\item
  ✅ MLRoute 智能路由（LightGBM, 95.8\%）
\item
  ✅ BudgetTracker / ConvergenceTracker
\item
  ✅ 549 测试，99.8\% 通过
\item
  ✅ Goedel-Prover-V2-8B 本地推理（vLLM, 24.38 tok/s）
\item
  ✅ DeepSeek V4 Flash / Pro API 集成
\item
  ✅ \textbf{社区物理库就绪}：Physlib（通用物理）+
  QuantumInfo（量子信息），Lean v4.30.0，3,169 commits
\end{itemize}

\section{待完成}\label{sec-todo}

\begin{itemize}
\tightlist
\item
  ⬜ Goedel-Prover-V2 全量 MiniF2F T2 验证
\item
  ⬜ 自动形式化（NL→Lean4）+ 验证管线
\item
  ⬜ \textbf{Physlib 物理定理自动证明} --- 在已有 Physlib/QuantumInfo
  定义基础上运行 Omega 证明引擎
\item
  ⬜ \textbf{QuantumInfo 量子光学扩展} ---
  连续变量体系、相干态、压缩态的 Physlib 贡献
\item
  ⬜ 持续学习：从失败证明中积累策略
\item
  ⬜ SorryDB 对接
\end{itemize}

\section{路线图}\label{sec-roadmap}

\subsection{短期（2026 Q3）}\label{ux77edux671f2026-q3}

{\def\LTcaptype{none} % do not increment counter
\begin{longtable}[]{@{}
  >{\raggedright\arraybackslash}p{(\linewidth - 2\tabcolsep) * \real{0.4286}}
  >{\raggedright\arraybackslash}p{(\linewidth - 2\tabcolsep) * \real{0.5714}}@{}}
\toprule\noalign{}
\begin{minipage}[b]{\linewidth}\raggedright
任务
\end{minipage} & \begin{minipage}[b]{\linewidth}\raggedright
里程碑
\end{minipage} \\
\midrule\noalign{}
\endhead
\bottomrule\noalign{}
\endlastfoot
全量 MiniF2F T2 验证运行 & 启动全量 488 题多轮证明 \\
Goedel-Prover-V2 本地自修正循环启用 & 编译错误反馈闭环 \\
SorryDB 对接与评估 & 接入开源证明任务基准 \\
Mathesis 自动形式化集成 & NL→Lean4 翻译管线 \\
\textbf{Physlib 集成验证} & 选取 Physlib 中简单物理定理，运行 Omega
自动证明 \\
\end{longtable}
}

\subsection{中期（2026 Q4）}\label{ux4e2dux671f2026-q4}

{\def\LTcaptype{none} % do not increment counter
\begin{longtable}[]{@{}
  >{\raggedright\arraybackslash}p{(\linewidth - 2\tabcolsep) * \real{0.4286}}
  >{\raggedright\arraybackslash}p{(\linewidth - 2\tabcolsep) * \real{0.5714}}@{}}
\toprule\noalign{}
\begin{minipage}[b]{\linewidth}\raggedright
任务
\end{minipage} & \begin{minipage}[b]{\linewidth}\raggedright
里程碑
\end{minipage} \\
\midrule\noalign{}
\endhead
\bottomrule\noalign{}
\endlastfoot
\textbf{Physlib 物理定理自动证明 demo} & 在 Physlib/ClassicalMechanics
上运行 Omega，验证端到端管线 \\
\textbf{基于 QuantumInfo 的量子光学自动证明} &
\([\hat{a}, \hat{a}^\dagger] = 1\), 密度矩阵迹运算等基础定理自动证明 \\
\textbf{连续变量体系形式化} & 相干态、压缩态的 Physlib 扩展定义 + Omega
证明 \\
持续学习管线 & 失败证明→策略积累→模型更新 \\
\end{longtable}
}

\subsection{长期（2027）}\label{ux957fux671f2027}

{\def\LTcaptype{none} % do not increment counter
\begin{longtable}[]{@{}
  >{\raggedright\arraybackslash}p{(\linewidth - 2\tabcolsep) * \real{0.5000}}
  >{\raggedright\arraybackslash}p{(\linewidth - 2\tabcolsep) * \real{0.5000}}@{}}
\toprule\noalign{}
\begin{minipage}[b]{\linewidth}\raggedright
任务
\end{minipage} & \begin{minipage}[b]{\linewidth}\raggedright
目标
\end{minipage} \\
\midrule\noalign{}
\endhead
\bottomrule\noalign{}
\endlastfoot
\textbf{Physlib 社区贡献管线} & Omega 自动证明 → PR 到 Physlib 主库 →
审阅合并 \\
端到端论文公式验证管线 & 输入论文 PDF → 输出已验证 Lean 证明（基于
Physlib 定义） \\
IMO 级别全自动证明 & 含自形式化（NL→Lean4 无需人工干预） \\
物理形式化覆盖扩展 & 扩展到量子场论、广义相对论等 Physlib TODO
列表中的领域 \\
\end{longtable}
}

\newpage

\cleardoublepage
\phantomsection
\addcontentsline{toc}{part}{Appendices}
\appendix

\chapter{验证实例：Omega
端到端证明}\label{ux9a8cux8bc1ux5b9eux4f8bomega-ux7aefux5230ux7aefux8bc1ux660e}

\section{命题}\label{sec-proposition}

\textbf{定理（乘法零律）：} 对任意自然数 \(a\)，\(a \times 0 = 0\)

\[
\forall a \in \mathbb{N}:\ a \cdot 0 = 0
\]

在 Peano 算术中，乘法递归定义：

\begin{itemize}
\tightlist
\item
  \(0 \cdot n = 0\)
\item
  \(S(m) \cdot n = m \cdot n + n\)
\end{itemize}

\section{Omega 自动生成的 Lean 4 代码}\label{sec-lean-code}

Omega 通过 DeepSeek V4 JSON mode 自动生成以下代码：

\begin{Shaded}
\begin{Highlighting}[]
\NormalTok{import Mathlib}

\NormalTok{theorem mul\_zero (a : Nat) : a * 0 = 0 :=}
\NormalTok{by}
\NormalTok{  induction a with}
\NormalTok{  | zero =\textgreater{} rfl}
\NormalTok{  | succ a ih =\textgreater{}}
\NormalTok{    simp [Nat.succ\_mul, ih]}
\end{Highlighting}
\end{Shaded}

\subsection{逐行解析}\label{ux9010ux884cux89e3ux6790}

\textbf{第 1 行：定理声明}

\begin{Shaded}
\begin{Highlighting}[]
\NormalTok{theorem mul\_zero (a : Nat) : a * 0 = 0 :=}
\end{Highlighting}
\end{Shaded}

{\def\LTcaptype{none} % do not increment counter
\begin{longtable}[]{@{}ll@{}}
\toprule\noalign{}
语法成分 & 说明 \\
\midrule\noalign{}
\endhead
\bottomrule\noalign{}
\endlastfoot
\texttt{theorem} & 声明新定理（不可计算消去） \\
\texttt{mul\_zero} & 定理名 \\
\texttt{(a\ :\ Nat)} & 全称量化的自然数变量 \\
\texttt{a\ *\ 0\ =\ 0} & 命题：Nat 类型上的等式，类型为 Prop \\
\texttt{:=} & 定义分隔符，右侧为证明体 \\
\end{longtable}
}

\textbf{第 2-5 行：证明体}

\begin{Shaded}
\begin{Highlighting}[]
\NormalTok{by}
\NormalTok{  induction a with}
\end{Highlighting}
\end{Shaded}

{\def\LTcaptype{none} % do not increment counter
\begin{longtable}[]{@{}
  >{\raggedright\arraybackslash}p{(\linewidth - 2\tabcolsep) * \real{0.3529}}
  >{\raggedright\arraybackslash}p{(\linewidth - 2\tabcolsep) * \real{0.6471}}@{}}
\toprule\noalign{}
\begin{minipage}[b]{\linewidth}\raggedright
语法
\end{minipage} & \begin{minipage}[b]{\linewidth}\raggedright
类型论解释
\end{minipage} \\
\midrule\noalign{}
\endhead
\bottomrule\noalign{}
\endlastfoot
\texttt{by} & 进入 tactic 模式（而非提供显式 \(\lambda\) 项） \\
\texttt{induction\ a\ with} & 在 \(a\) 上归纳：基础情形 \(a=0\) +
归纳步骤 \(a=k+1\) \\
\end{longtable}
}

\textbf{第 3 行：基础情形}

\begin{Shaded}
\begin{Highlighting}[]
\NormalTok{  | zero =\textgreater{} rfl}
\end{Highlighting}
\end{Shaded}

{\def\LTcaptype{none} % do not increment counter
\begin{longtable}[]{@{}ll@{}}
\toprule\noalign{}
语法 & 说明 \\
\midrule\noalign{}
\endhead
\bottomrule\noalign{}
\endlastfoot
\texttt{\textbar{}\ zero\ =\textgreater{}} & 模式匹配：\(a=0\) 的情形 \\
\texttt{rfl} & 自反性：\(0 \cdot 0 = 0\) 在定义上相等 \\
\end{longtable}
}

\textbf{第 4-5 行：归纳步骤}

\begin{Shaded}
\begin{Highlighting}[]
\NormalTok{  | succ a ih =\textgreater{}}
\NormalTok{    simp [Nat.succ\_mul, ih]}
\end{Highlighting}
\end{Shaded}

{\def\LTcaptype{none} % do not increment counter
\begin{longtable}[]{@{}
  >{\raggedright\arraybackslash}p{(\linewidth - 2\tabcolsep) * \real{0.5000}}
  >{\raggedright\arraybackslash}p{(\linewidth - 2\tabcolsep) * \real{0.5000}}@{}}
\toprule\noalign{}
\begin{minipage}[b]{\linewidth}\raggedright
语法
\end{minipage} & \begin{minipage}[b]{\linewidth}\raggedright
说明
\end{minipage} \\
\midrule\noalign{}
\endhead
\bottomrule\noalign{}
\endlastfoot
\texttt{\textbar{}\ succ\ a\ ih\ =\textgreater{}} & 模式匹配：假设
\(a=k\) 时成立（归纳假设 \texttt{ih}），证明 \(a=k+1\) \\
\texttt{simp\ {[}Nat.succ\_mul,\ ih{]}} & 应用 Mathlib
引理和归纳假设化简 \\
\end{longtable}
}

\textbf{证明的逻辑链条：}

\[
\begin{aligned}
\text{base }(a=0)&: 0\cdot 0 = 0 \quad \text{(by definition of $\cdot$)} \\
\text{step }(a=k+1)&: (k+1)\cdot 0 = k\cdot 0 + 0 \quad \text{(by Nat.succ\_mul)} \\
&= 0 + 0 \quad \text{(by induction hypothesis $k\cdot 0 = 0$)} \\
&= 0 \quad \text{(by definition of $+$)}
\end{aligned}
\]

\section{编译结果}\label{sec-compile-result}

Omega 的输出记录：

\begin{verbatim}
=== E2E: Omega + DeepSeek v4 JSON + Lean4/Mathlib ===
generate_fn: OK
Theorem: theorem mul_zero (a : Nat) : a * 0 = 0 :=
JSON output:
  tactic: induction a with
          | zero => rfl
          | succ a ih => simp [Nat.succ_mul, ih]
  confidence: 0.9
  is_complete: true
  lean_code: (125 chars, with import Mathlib)

=== T2 COMPILE ===
exit_code: 0
stderr: (empty)
stdout: (empty)

✅ END-TO-END PASS
\end{verbatim}

{\def\LTcaptype{none} % do not increment counter
\begin{longtable}[]{@{}ll@{}}
\toprule\noalign{}
关键指标 & 数值 \\
\midrule\noalign{}
\endhead
\bottomrule\noalign{}
\endlastfoot
模型 & DeepSeek V4 Flash（JSON mode） \\
推理耗时 & \textasciitilde2s \\
JSON 置信度 & 0.9（模型自报） \\
T2 编译耗时 & \textasciitilde2.5s（含 Mathlib 加载） \\
最终状态 & exit\_code=0，编译通过 \\
\end{longtable}
}

该证明使用归纳法，正确调用了 Mathlib 引理
\texttt{Nat.succ\_mul}。\texttt{rfl}
处理基础情形，\texttt{simp\ {[}Nat.succ\_mul,\ ih{]}}
处理归纳步，代码整洁、结构性完整。

\newpage

\chapter{MiniF2F 涵盖数据集}\label{sec-dataset-detail}

{\def\LTcaptype{none} % do not increment counter
\begin{longtable}[]{@{}lll@{}}
\toprule\noalign{}
子集 & 题目数 & 来源 \\
\midrule\noalign{}
\endhead
\bottomrule\noalign{}
\endlastfoot
\texttt{mathd\_algebra} & 56 & MATH 数据集代数子集 \\
\texttt{mathd\_numbertheory} & 30 & MATH 数据集数论子集 \\
\texttt{mathd\_prealgebra} & 42 & MATH 数据集预代数子集 \\
\texttt{mathd\_precalculus} & 14 & MATH 数据集预微积分子集 \\
\texttt{amc12} & 26 & AMC 12 竞赛题 \\
\texttt{aime} & 50 & AIME 竞赛题 \\
\texttt{imo} & 8 & IMO 竞赛题 \\
\texttt{induction} & 44 & 归纳法专项 \\
\texttt{algebra} & 76 & 代数综合 \\
\texttt{numbertheory} & 40 & 数论综合 \\
\textbf{MiniF2F 总计} & \textbf{488} & --- \\
\end{longtable}
}

\newpage

\chapter{技术栈依赖}\label{ux6280ux672fux6808ux4f9dux8d56}

{\def\LTcaptype{none} % do not increment counter
\begin{longtable}[]{@{}lll@{}}
\toprule\noalign{}
组件 & 版本 & 用途 \\
\midrule\noalign{}
\endhead
\bottomrule\noalign{}
\endlastfoot
Lean 4 & 4.30.0 & 定理证明器 \\
Mathlib & \textasciitilde2.0 (8109 oleans, 6.7 GB) & 数学定理库 \\
Python & 3.13 & 系统语言 \\
vLLM & 0.22.1 & Goedel-Prover 推理引擎 \\
DeepSeek API & --- & 远程 LLM 推理 \\
LightGBM & --- & ML 路由分类器 \\
ONNX Runtime & --- & ML 路由推理 \\
json-repair & 0.60.1 & JSON 容错修复 \\
\end{longtable}
}

\newpage

\chapter{元数据标准与依赖图}\label{ux5143ux6570ux636eux6807ux51c6ux4e0eux4f9dux8d56ux56fe}

\section{章节元数据字段}\label{ux7ae0ux8282ux5143ux6570ux636eux5b57ux6bb5}

{\def\LTcaptype{none} % do not increment counter
\begin{longtable}[]{@{}llll@{}}
\toprule\noalign{}
字段 & 类型 & 说明 & 示例 \\
\midrule\noalign{}
\endhead
\bottomrule\noalign{}
\endlastfoot
\texttt{chapter-id} & str & 唯一标识 & \texttt{ch03} \\
\texttt{chapter-version} & str & 章节版本 & \texttt{0.1.0} \\
\texttt{category} & str & 分类 & \texttt{usage} \\
\texttt{audience} & list{[}str{]} & 目标读者 &
\texttt{{[}AI4Math算法组,\ 系统工程师{]}} \\
\texttt{keywords} & list{[}str{]} & 关键词 &
\texttt{{[}Quick\ Start,\ 安装{]}} \\
\texttt{git-commit} & str & 提交哈希 & (脚本注入) \\
\texttt{last-reviewed} & date & 审阅日期 & \texttt{2026-06-10} \\
\texttt{prerequisites} & list{[}str{]} & 前置知识 &
\texttt{{[}系统架构{]}} \\
\texttt{depends-on} & list{[}str{]} & 依赖章节 & \texttt{{[}ch02{]}} \\
\end{longtable}
}

\section{章节依赖图}\label{ux7ae0ux8282ux4f9dux8d56ux56fe}

\begin{forest}
for tree={grow=south,edge={gray!60},l sep=8pt,
  font=\small,parent anchor=north,child anchor=south}
[Ω-Architect 用户指南
  [idx: 索引页]
  [ch01: 项目概述
    [ch02: 系统架构
      [ch03: 性能指标]
      [ch04: Quick Start]
    ]
  ]
  [ch05: 生态与对比
    [ch06: 路线图]
  ]
  [appx: 验证实例 \& 数据]
]
\end{forest}

\section{版本锚点注入}\label{ux7248ux672cux951aux70b9ux6ce8ux5165}

\texttt{0.1.0}, \texttt{2953191}, \texttt{2026-06-10} 由
\texttt{scripts/generate\_user\_guide\_metadata.py} 自动替换。

\textbf{规则：} 用户指南版本与 \texttt{pyproject.toml} 中的 repo
版本同步。PDF 渲染前注入 HEAD commit hash 和日期。渲染后
\texttt{-\/-restore} 恢复占位符，源文件保持 git 干净。

\chapter{版本历史}\label{ux7248ux672cux5386ux53f2}

{\def\LTcaptype{none} % do not increment counter
\begin{longtable}[]{@{}lllll@{}}
\toprule\noalign{}
用户指南版本 & repo 版本 & repo commit & 日期 & 变更 \\
\midrule\noalign{}
\endhead
\bottomrule\noalign{}
\endlastfoot
v0.1.0 & 0.1.0 & \texttt{2953191} & 2026-06-10 & 初始用户指南框架 \\
\end{longtable}
}

\textbf{版本同步规则：}

\begin{enumerate}
\def\labelenumi{\arabic{enumi}.}
\tightlist
\item
  \textbf{repo 版本} 在 \texttt{pyproject.toml} 的
  \texttt{{[}project{]}\ version} 中维护，是唯一数据源
\item
  \textbf{文档版本} 与 repo 版本保持一致，在 \texttt{index.qmd} 和每章
  \texttt{doc-version} 字段中标注
\item
  \textbf{git commit hash} 在 PDF 渲染前注入，提供精确的可溯源指针
\item
  \textbf{渲染流程：}
  \texttt{python3\ scripts/generate\_user\_guide\_metadata.py} →
  \texttt{quarto\ render} → \texttt{-\/-restore}
\end{enumerate}




\end{document}
